\par\medskip\noindent\textbf{LLM single call, LLM monotone, LLM monotone (constrained) and LLM re-rank: system message.}\par\nobreak
\begin{lstlisting}
You align lecture transcripts to slide decks. You answer with JSON and nothing else: no explanation, no code fence, no prose before or after the object.
\end{lstlisting}

\par\medskip\noindent\textbf{LLM single call (also the LLM voter of Ensemble vote): user message.}\par\nobreak
\begin{lstlisting}
The deck has {n_slides} slides, numbered 0 to {last_slide}.
The transcript has {n_rows} rows, numbered 0 to {last_row}.

SLIDES
{slides}

TRANSCRIPT
{rows}

Assign every transcript row to the slide it is about. Rows about the same slide are
usually next to each other, but you are not required to move through the deck in order,
and slides nobody discussed may be left out.

Answer with a JSON object of exactly this shape:

  {"blocks": [{"from": a, "to": b, "slide": j}, {"from": c, "to": d, "slide": k}, ...]}

"from" and "to" are INCLUSIVE row numbers, "slide" is a slide number in
0..{last_slide}, and the blocks together must cover every row from 0 to {last_row}
exactly once, listed in increasing order of "from".
\end{lstlisting}

\par\medskip\noindent\textbf{LLM monotone: user message.}\par\nobreak
\begin{lstlisting}
The deck has {n_slides} slides, numbered 0 to {last_slide}.
The transcript has {n_rows} rows, numbered 0 to {last_row}.

SLIDES
{slides}

TRANSCRIPT
{rows}

The lecturer moves through the deck in order, from slide 0 to slide {last_slide},
without going back. Say where each slide STARTS: the row number at which the lecturer
begins talking about it.

Answer with a JSON object of exactly this shape:

  {"starts": [0, s1, s2, ...]}

The list has exactly {n_slides} numbers, one per slide in deck order. It must be
non-decreasing, the first number must be 0, and every number must be in
0..{last_row}. A slide the lecturer skipped gets the same start as the slide after it.
\end{lstlisting}

\par\medskip\noindent\textbf{LLM monotone (constrained): user message.}\par\nobreak
\begin{lstlisting}
The deck has {n_slides} slides, numbered 0 to {last_slide}.
The transcript has {n_rows} rows, numbered 0 to {last_row}.

SLIDES
{slides}

TRANSCRIPT
{rows}

The lecturer moves through the deck in order, from slide 0 to slide {last_slide},
without ever going back. Say how LONG the lecturer spends on each slide: for each slide
in deck order, the number of consecutive transcript rows spent on it.

Answer with a JSON object of exactly this shape:

  {"gaps": [g0, g1, g2, ...]}

The list has exactly {n_slides} numbers, one per slide in deck order, each of them 0 or
more. A slide the lecturer skipped gets 0. The numbers should add up to about {n_rows},
the length of the transcript.

Slide 0 starts at row 0. Slide k starts at the row given by the sum of the first k
numbers, so the alignment always moves forward through the deck.
\end{lstlisting}

\par\medskip\noindent\textbf{LLM re-rank (5 methods): user message.}\par\nobreak
\begin{lstlisting}
The deck has {n_slides} slides, numbered 0 to {last_slide}.
The transcript has {n_rows} rows, numbered 0 to {last_row}.

SLIDES
{slides}

TRANSCRIPT
{rows}

Below are {k} candidate alignments of this transcript to this deck. Each was produced
by a different alignment method and is written as the row ranges on which each slide is
being discussed (a slide may appear more than once if the lecturer returns to it).

{candidates}

Exactly one of them is to be chosen. Pick the candidate whose slide ranges best match
what the transcript actually says, and answer with a JSON object of exactly this shape:

  {"choice": 1}

where the number is between 1 and {k}.
\end{lstlisting}

\par\medskip\noindent\textbf{LLM hierarchical (also inside LLM-prior DP), split call: system message.}\par\nobreak
\begin{lstlisting}
You are aligning lecture transcript packets to slide-page ranges.

The slides and transcript are from the SAME lecture.
The spoken transcript does NOT need exact wording overlap with the slide text.
A transcript packet can belong to a page range even when it:
- paraphrases slide content
- expands with examples
- explains figures/tables indirectly
- continues the instructor's explanation without repeating slide words

You are given:
- a LEFT consecutive page range
- a RIGHT consecutive page range
- consecutive transcript packets from one continuous lecture segment

Your job:
decide where the lecture moves from the LEFT page range to the RIGHT page range.

Return strict JSON only.
\end{lstlisting}

\par\medskip\noindent\textbf{LLM hierarchical, split call: user message.}\par\nobreak
\begin{lstlisting}
LEFT PAGE RANGE:
{left_pages}

RIGHT PAGE RANGE:
{right_pages}

TRANSCRIPT PACKETS IN CHRONOLOGICAL ORDER:
{packets}

Return strict JSON with this exact schema:
{
  "transition_after_local_packet_pos": <int or null>
}

Interpretation:
- if transition_after_local_packet_pos is X, packets up to local position X belong mainly to the LEFT page range,
  and packets after X belong mainly to the RIGHT page range
- local packet positions are 0-based within the packet list shown above
- if ALL packets belong mainly to the RIGHT page range, return null
- if ALL packets belong mainly to the LEFT page range, return the last local packet position shown

Rules:
- preserve chronology
- be conservative
- choose the most plausible transition in lecture flow, not literal text overlap
- do not explain your choice and do not output any text outside the JSON object

Pacing guidance: on average about {spp} sections correspond to one slide, so a page range of P pages should receive roughly P*{spp} sections. Prefer a transition that reflects this pacing over dumping all packets on one side.
\end{lstlisting}

\par\medskip\noindent\textbf{LLM hierarchical, leaf call: system message.}\par\nobreak
\begin{lstlisting}
You are aligning transcript sections to a small consecutive range of slide pages.

The slides and transcript are from the SAME lecture.
A section may belong to a page even if the wording differs. It may paraphrase, expand, explain examples, or continue the instructor's explanation of the page.

Return a monotonic section-to-page assignment:
- section assignments must be non-decreasing in page index
- pages may receive zero or more sections
- use exact provided page indices only
Return strict JSON only.
\end{lstlisting}

\par\medskip\noindent\textbf{LLM hierarchical, leaf call: user message.}\par\nobreak
\begin{lstlisting}
PAGES:
{pages}

TRANSCRIPT SECTIONS:
{sections}

Return strict JSON with this exact schema:
{
  "assignments": [
    {"section_idx": <int>, "page_index": <int>}
  ]
}

Rules:
- assign every shown section exactly one page_index from the provided pages
- assignments must be monotonic non-decreasing in page_index
- use lecture-flow reasoning, not literal text matching only
- do not give reasons and do not output any text outside the JSON object

Pacing guidance: in this lecture the instructor spends on average about {spp} transcript sections per slide. Distribute the sections across the pages accordingly; do NOT lump most sections onto the first page. Move to the next page once the discussion clearly shifts to it, and keep progressing in order.
\end{lstlisting}
